\documentclass[11pt]{letter}
\usepackage[margin=1in]{geometry}
\usepackage{hyperref}
\usepackage{parskip}

\signature{Christopher Regan \\ Ying Xie}
\address{College of Computing and Software Engineering \\ Kennesaw State University \\ Marietta, GA 30060, USA \\ cregan1@kennesaw.edu}

\begin{document}

\begin{letter}{The Editors \\ Digital Finance \\ Springer}

\opening{Dear Editors,}

We are pleased to submit our manuscript, ``Validating LLM Structural Reasoning: Detecting Persistent Market Regimes Through Temporal Obfuscation,'' for consideration as an Original Research Article in \textit{Digital Finance}.

\textbf{Context and importance.}
A central challenge in deploying large language models for financial analysis is distinguishing genuine structural reasoning from training data memorization. We introduce \textit{temporal obfuscation testing}---systematically stripping dates, ticker symbols, and contextual markers from input sequences---as a general-purpose methodology to address this challenge. Applied to options dealer gamma exposure (GEX) data spanning 2020--2025, the framework validates LLM regime detection across 2,221 evaluations (1,412 real windows and 809 synthetic controls), achieving 81.2\% detection of persistent dealer regimes in 2024 versus 12.1\% in 2020, with 0\% false positives on negative controls.

\textbf{Key contributions.}
The manuscript offers three results we believe are of particular interest to the \textit{Digital Finance} readership:

\begin{enumerate}
\item \textbf{Raw data superiority over parametric compression.} When provided raw strike-level options chain data instead of pre-calculated GEX metrics, the LLM achieves 92.3\% detection versus 61.5\%---a 30.8 percentage point improvement. This finding challenges standard data pipeline designs that aggregate information before analysis and demonstrates that scalar metrics represent lossy compression of structural signal.

\item \textbf{0DTE-driven market structure evolution.} Multi-year analysis reveals gradual regime evolution tracking zero-days-to-expiration options adoption, with detection progressing from 3.7\% (2021) to 100\% (2024) and GEX magnitude growing 360\%. This documents a structural reorganization of dealer positioning with direct implications for market surveillance and risk management.

\item \textbf{Detection--alpha orthogonality.} Stable pattern detection (68--74\% quarterly) persists as trading profitability collapses (Sharpe ratio 1.8 to 0.1), establishing that detected patterns represent structural market mechanics rather than exploitable inefficiencies---an important distinction for responsible AI deployment in finance.
\end{enumerate}

\textbf{Fit with Digital Finance.}
The manuscript sits at the intersection of AI methodology validation and financial market microstructure---core areas of the journal's scope. The temporal obfuscation framework generalizes beyond finance, but its application to options dealer positioning, 0DTE market structure, and the information-theoretic implications of data granularity for AI-driven analysis addresses questions directly relevant to the digital transformation of financial markets.

\textbf{Disclosures.}
This manuscript has not been published elsewhere and is not under consideration by another journal. A companion single-day study was published at IEEE BigData 2025 (arXiv:2512.17923); the present work substantially extends that study with multi-day regime detection, negative controls, multi-year temporal analysis, and theoretical framing. All authors have approved the manuscript and agree to its submission. The authors declare no conflicts of interest. All code and data are publicly available at \url{https://github.com/iAmGiG/gex-llm-patterns}.

We believe this work will be of interest to \textit{Digital Finance} readers working at the intersection of AI validation, market microstructure, and financial data science. We look forward to your consideration.

\closing{Sincerely,}

\end{letter}
\end{document}
